% Section 7: Discussion
% Source: context/Discussion_and_Conclusions.md §D2-D8
% Source: context/ICSA_RA_Literature_Review.md §12
%
% EM-DASH PROHIBITION: Do not use --- or \textemdash. Use colons or semicolons.

\section{Discussion}\label{sec:discussion}

The empirical findings and case study evidence converge on a single conclusion:
LLM credential relay corruption is a systematic failure mode that prompt
engineering did not eliminate in any condition tested, motivating an architectural
response. We discuss six implications.

\textbf{Prompt engineering as insufficient mitigation.}
The production system deployed the most explicit bearerToken instruction
achievable within the MCP tool schema (318 characters naming the RSA mechanism,
the 403 consequence, and single-character precision). It produced 19.9\%
\DFR{} ($p\!=\!0.16$; a null result, not a positive equivalence claim). BPE tokenisation
boundaries \emph{appear} to operate below the instruction-following layer; no
instruction eliminated corruption in any of the 16 models tested. In the conditions
evaluated, prompt improvement was bounded at $\sim$5~pp (L1-to-L3) with no approach toward zero. We present
this as an \emph{observed practical insufficiency}: the most aggressive prompting
strategy available under MCP schema constraints did not reduce \DFR{} to an
acceptable level for cryptographic relay. Formal equivalence testing (TOST with a pre-specified margin)
remains future work; the current finding is sufficient to motivate an
architectural bypass.

\textbf{Mitigation vs.\ bypass: a category difference.}
The Vault achieves zero \DFR{} by construction for the byte-relay property
(a design argument, not an additional empirical measurement): removing the
credential from inference makes byte-level corruption architecturally
impossible. This guarantee is scoped to relay fidelity; it does not address
credential theft at rest, token expiry, revocation, or replay attacks;
those remain the responsibility of the surrounding identity infrastructure. Prompt-based mitigation asks how to
make the LLM relay more reliably; \DFR{} remains above zero. These are
solutions to different problems. Bypass is appropriate when correctness is
required (RSA verification is binary), when the failure is
unrecoverable (HTTP~403 with no retry), and when credentials are sensitive personal data~\cite{gdpr2016}.

\textbf{When to apply the Vault.}
Apply when: (i)~credential length exceeds $\sim$200~chars;
(ii)~a cryptographic signature is present; or (iii)~a relay error is unrecoverable.
Prompt mitigation may suffice in retry-tolerant contexts below that threshold.

\textbf{Organic convergence validates the RA.}
The production Vault was implemented independently without reference to
this research (Section~\ref{sec:trajectory}), confirming the RA as
naturally emergent and transferable without requiring rediscovery.

\textbf{Credential lengths are increasing; the risk may worsen.}
Gibberish bias intensifies with larger vocabularies~\cite{chen2026gibberish}.
NIST key sizes~\cite{barker2020nistkeys} and post-quantum
algorithms~\cite{nist2024pqc} (CRYSTALS-Dilithium: 2{,}420~bytes; SPHINCS+:
7{,}856~bytes) would push tokens deeper into the high-\DFR{} tier, making 17.3\% a
conservative lower bound (this projection requires empirical validation). The
Vault addresses this trajectory without modification.

\textbf{Broader implications for agentic AI systems.}
As LLM-based agent systems proliferate~\cite{wang2024survey}, any opaque
high-entropy string (API keys, hash digests, certificate fragments) may share
the same distribution-shift vulnerability. JWT relay is the validated case;
extension to other credential classes is a hypothesis requiring independent
empirical confirmation. The eco arm
($\rho\!=\!0.740$) provides a reusable calibration instrument for
deployment-specific risk.
